\section{Experiments}\label{sec:experiments}
\subsection{Experimental setup}
\paragraph{Data and metrics}
Identify dataset versions, licenses, splits, exclusions, and metrics. Report
the unit of aggregation and whether intervals describe variation across
examples, seeds, or independent runs. State how overlap and contamination
were checked.

\paragraph{Implementation details}
Specify model checkpoints, software versions, hyperparameters, random seeds,
hardware, and the number of optimization steps. Report the actual resource
configuration used by each experiment.

\subsection{Results and ablations}
Table~\ref{tab:results} uses a full-width three-rule table; values are left
blank intentionally. Table captions appear above the table, and figure
captions appear below the figure.
\begin{table}[ht]
\caption{Main comparison. Dashes indicate unfilled placeholders. No values in
this template are experimental measurements.}\label{tab:results}
\centering
\begin{tabular*}{\linewidth}{@{\extracolsep{\fill}}llccc@{}}
\toprule
Method & Setting & Quality $\uparrow$ & Latency $\downarrow$ & Memory $\downarrow$\\
\midrule
Baseline & Reference & -- & -- & --\\
Proposed & Full method & -- & -- & --\\
\midrule
Ablation & Without component A & -- & -- & --\\
Ablation & Without component B & -- & -- & --\\
\bottomrule
\end{tabular*}
\end{table}

Discuss accuracy, efficiency, and failure modes together. Do not attribute a
gain to a component without an appropriate control experiment.

\begin{figure}[ht]
\centering
\begin{subfigure}[t]{0.47\linewidth}
  \centering
  \fbox{\parbox[c][0.85in][c]{0.88\linewidth}{\centering Qualitative example A}}
  \caption{Describe the input and prediction.}\label{fig:case-a}
\end{subfigure}\hfill
\begin{subfigure}[t]{0.47\linewidth}
  \centering
  \fbox{\parbox[c][0.85in][c]{0.88\linewidth}{\centering Qualitative example B}}
  \caption{Describe a representative failure.}\label{fig:case-b}
\end{subfigure}
\caption{Qualitative examples. Use the same crop, label size, and color
conventions across panels. These boxes are placeholders.}\label{fig:cases}
\end{figure}

\subsection{Code and reproducibility}
Listing~\ref{lst:config} demonstrates line-numbered code without requiring
shell escape or external syntax-highlighting programs.
\begin{lstlisting}[language=Python,caption={Illustrative evaluation configuration.},label={lst:config}]
config = {
    "seed": 42,
    "split": "validation",
    "batch_size": 1,
    "checkpoint": "REPLACE_WITH_CHECKPOINT",
}
# Record exact versions and commands with each run.
\end{lstlisting}
For long URLs, use \verb|\url{...}| rather than manually inserting line breaks.
Use footnotes sparingly.\footnote{An example footnote, set smaller than the
body text and placed at the bottom of the relevant page.}
